\chapter{Анализ предметной области}\label{ch:analysis}

\section{Постановка задачи}\label{sec:analysis-problem}

В этом разделе зафиксируйте исходную проблему, ограничения и ожидаемый
результат магистерской работы. Удобно отдельно описать входные данные, условия
применимости и критерии успешности решения.

\section{Обзор существующих подходов}\label{sec:analysis-review}

Здесь обычно приводят сравнение методов, инструментов, алгоритмов или
технологий, на которые опирается работа. Для каждого варианта полезно кратко
указать преимущества, ограничения и причину, по которой он выбран или
отклонён.

\subsection{Критерии сравнения}\label{subsec:analysis-criteria}

Сформулируйте критерии, по которым оцениваются рассматриваемые подходы:
точность, скорость, стоимость реализации, удобство сопровождения, требования к
данным и другие релевантные параметры.

\subsection{Выводы по обзору}\label{subsec:analysis-summary}

Подведите итог аналитической части и обоснуйте, какое направление будет
развиваться далее в проектной и практической части работы.
